% Clase 20 — tambor estacionario con un taco. Esquema cualitativo, no una
% solución calculada. Izquierda: vista de arriba; el borde está a altura cero y
% la zona del taco (el docente dibuja un hexágono) a altura h_0; las curvas
% punteadas sugieren curvas de nivel. Derecha: corte por el diámetro; fuera del
% taco la altura decae hasta el borde cumpliendo Laplace.
\begin{tikzpicture}[scale=1]

  % ---------- vista de arriba ----------
  \fill[figgray!10] (0,0) circle (2.2);
  \draw[figline, line width=1.1pt] (0,0) circle (2.2);
  \draw[figaux] (0.25,0.1) ellipse (1.65 and 1.6);
  \draw[figaux] (0.4,0.15) ellipse (1.15 and 1.1);
  \fill[figamber!30, draw=figamber, line width=0.9pt]
    (1.2,0.2) -- (0.85,0.806) -- (0.15,0.806) -- (-0.2,0.2) -- (0.15,-0.406) -- (0.85,-0.406) -- cycle;
  \node[figlbl] at (0.5,0.52) {$h_0$};
  \node[figlbl] at (-1.75,-1.95) {$h = 0$};
  \draw[figgray, dash dot, line width=0.5pt] (-2.5,0.2) -- (2.5,0.2);
  \node[figlbl, anchor=north] at (0,-2.5) {(a) vista de arriba};

  % ---------- corte ----------
  \begin{scope}[xshift=4.0cm, yshift=-0.6cm]
    \draw[figaxis] (-0.2,0) -- (5.1,0);
    \node[figlbl] at (5.35,0) {$x$};
    \draw[figaxis] (0,-0.2) -- (0,2.3);
    \node[figlbl] at (0,2.55) {$h$};
    % lonja
    \draw[figblue, line width=1.2pt]
      (0.3,0) .. controls (1.4,0.15) and (1.9,1.5) .. (2.3,1.5) -- (3.7,1.5) .. controls (3.95,1.5) and (4.3,0.25) .. (4.7,0);
    % taco
    \fill[figamber!30, draw=figamber, line width=0.9pt] (2.3,0) rectangle (3.7,1.5);
    \node[figlbl] at (3.0,0.75) {taco};
    \draw[figaux] (0,1.5) -- (2.3,1.5);
    \node[figlblaux, anchor=east] at (-0.08,1.5) {$h_0$};
    % borde
    \fill[figgray] (0.3,0) circle (1.6pt);
    \fill[figgray] (4.7,0) circle (1.6pt);
    \node[figlblaux] at (0.3,-0.3) {borde};
    \node[figlblaux] at (4.7,-0.3) {borde};
    \node[figlbl, anchor=north] at (2.5,-1.9) {(b) corte por la línea de trazos};
  \end{scope}

\end{tikzpicture}
